%
% $Id: slides.tex 4228 2006-06-21 21:55:12Z jjamor $
%
%
% Compilar a .pdf con LaTeX (pdflatex)
% Es necesario instalar Beamer (paquete latex-beamer en Debian)
%

%
% Gráficos:
% Los gráficos pueden suministrarse en PNG, JPG, TIF, PDF, MPS
% Los EPS deben convertirse a PDF (usar epstopdf)
%

\documentclass{beamer}
\usetheme{Warsaw}
\usebackgroundtemplate{\includegraphics[width=\paperwidth]{format/libresoft-bg.png}}
\usepackage[spanish]{babel}
\usepackage[utf8]{inputenc}
\usepackage{graphics}
\usepackage{amssymb} % Simbolos matematicos
\usepackage{url}

%\definecolor{libresoftgreen}{RGB}{162,190,43}
%\definecolor{libresoftblue}{RGB}{0,98,143}

%\setbeamercolor{titlelike}{bg=libresoftgreen}

%% Metadatos del PDF.
\hypersetup{
  pdftitle={Debate open innovation video (C. Leadbeater)},
  pdfauthor={F. Ortega, Jesus M. Gonzalez-Barahona},
  pdfcreator={GSyC/Libresoft},
  pdfproducer=PDFLaTeX,
  pdfsubject={nn},
}
%%


\AtBeginSection[]
{
  \begin{frame}<presentation>
    \frametitle{Index}
    \tableofcontents[current]
  \end{frame}
}


\begin{document}

\title{Debate open innovation video (C. Leadbeater)}
\subtitle{Master on Free Software}
\institute{\\jfelipe@libresoft.es\\
GSyC/Libresoft}
\author{Felipe Ortega, Jesus M. Gonzalez-Barahona}
\date{\today}

\frame{
\maketitle
\begin{center}
\includegraphics[width=6cm]{format/gsyc-urjc}
\end{center}
}


% Si el titulo o el autor se quieren acortar para los pies de p�gina
% se pueden redefinir aqu�:
%\title{Titulo corto}
%\author{Autores abreviado}


%% LICENCIA DE REDISTRIBUCION DE LAS TRANSPAS
\frame{
~
\vspace{4cm}

\begin{flushright}
{\tiny
(cc) 2010 Felipe Ortega, Jesus M. Gonzalez-Barahona. \\
Some rights reserved. This document is distributed under the Creative \\
            Commons Attribution-ShareAlike 3.0 licence, available in \\
            http://creativecommons.org/licenses/by-sa/3.0/

%  Este documento (o uno muy similar) está disponible en \\
%  \url{http://gsyc.escet.urjc.es/~jjamor/}
}
\end{flushright}
}
%%

%%%%%%
%Transpas separadas por \begin{frame}
%%%%%%%%%%%%%%%%%%%%%%%%\end{frame}

%\section{Aim of this exercise}

\begin{frame}
\frametitle{Dynamics}
\begin{itemize}
\item Some references from video presentation.
\item Possible topics or questions for open debate.
\end{itemize}
\end{frame}

%%%%%%%%%%%%%%%%%%%%%%%%%%%%%%%%%%%%%%%%%%%%%%%%%%%%%%%%%%%%%%

\begin{frame}
 \frametitle{Internet changes information economy}
 \begin{itemize}
  \item 19-20th century: High cost as initial requirement for starting mass circulation of info.
  \item \textit{http://setiathome.berkeley.edu/}
  \item \textit{http://www.chem.ox.ac.uk/cancer/pancreaticcancer.html}
 \end{itemize}


\end{frame}

%%%%%%%%%%%%%%%%%%%%%%%%%%%%%%%%%%%%%%%%%%%%%%%%%%%%%%%%%%%%%%

\begin{frame}
 \frametitle{Needed capital is distributed}
 \begin{itemize}
  \item Apache vs. Microsoft IIS.
  \item \textit{http://beamartian.jpl.nasa.gov/welcome}
  \item \textit{Crowdsourcing}, \textit{wisdom of crowds}...
 \end{itemize}

\end{frame}

%%%%%%%%%%%%%%%%%%%%%%%%%%%%%%%%%%%%%%%%%%%%%%%%%%%%%%%%%%%%%%

\begin{frame}
 \frametitle{New production models, massive distributed communities}
 \begin{itemize}
  \item Google PageRank.
  \item Wikipedia vs. Encarta.
  \item \textit{http://www.crowdspring.com/}
  \item \textit{P2P networks}
 \end{itemize}

\end{frame}

\end{document}